\chapter{Cenni storici su TCFS}
Nella sua versione originale TCFS era molto diverso diverso da quello che è oggi, tutte le seguenti considerazioni in questo capitolo sono frutto dell'analisi di un unico documento, una delle ultime tracce rimaste della versione originale di TCFS. Il documento denominato: "The Design and Implementation of a Transparent Cryptographic Filesystem for UNIX" presentato nel 2001 allo USENIX è stato di forte ispirazione per il progetto. Sfortunatamente il codice originale è stato perso, questo ha comportato una ri-scrittura totale del FS.

\section{Il vecchio TCFS: origine e caratteristiche}
TCFS inizia a prendere forma ad UNISA, nel 1995, con lo scopo di rendere sicura la conservazione dei dati su macchine un pò diverse da quelle che abbiamo oggi, si intende dire che i computer erano basati su un’architettura in cui c’era una macchina principale e diversi terminali connessi. Questo era dovuto all’alto costo delle macchine dell’epoca, che rendeva impossibile avere un computer prestante e personale per ogni utente. Ad oggi diamo per scontato che il computer sia personale e che di solito l'utente connesso a una macchina sia singolo. Nella nuova versione di TCFS si sono tenuti in considerazione questi aspetti.

\subsection{Il contesto storico e le motivazioni del progetto}
La protezione dei dati memorizzati sul server era ed è tutt'ora una sfida importante, in quanto il server può essere soggetto a attacchi, guasti, o accessi non autorizzati. Inoltre, la comunicazione tra il client e il server può essere intercettata o alterata da terze parti malintenzionate. Per questi motivi si sentiva la necessità di sviluppare dei filesystem crittografici che potessero garantire la confidenzialità, l’integrità, e l’autenticazione dei dati, sia in transito che in riposo. I filesystem crittografici sono stati oggetto di molte ricerche negli ultimi decenni, a partire dalle prime proposte. Lungo questa linea di ricerca, la trasparenza (il fatto che il filesystem non sia locale dovrebbe essere nascosto alle applicazioni e agli utenti) e la sicurezza sono state le due questioni principali.
Alcuni sistemi, come il Cluster File System di Compaq , garantiscono la trasparenza anche in presenza di guasti, ma assumono che la confidenzialità dei dati e l’autenticazione delle parti coinvolte sia ottenuta attraverso qualche altro sistema.
Il File System di Andrew e CODA invece forniscono dei meccanismi per autenticare i server su linee pubbliche e per assicurare che la comunicazione client-server non possa essere intercettata. Le serie questioni di gestione delle chiavi che sorgono in tali file system sono affrontate dal Self-certifying File System. Facciamo notare che sia AFS che CODA assumono che il server sia affidabile e memorizzano i dati in chiaro sulla macchina server. TCFS invece memorizza i file in forma criptata, negando al server l’accesso ai dati in chiaro.

\subsection{L’architettura e le funzionalità del vecchio TCFS}
Date le tecnologie dell'epoca TCFS era strutturado in maniera completamente diversa rispetto ad oggi. Esso era costituito da un modulo kernel scritto "ad hoc", che si occupava di memorizzare e mantenere in maniera sicura le chiavi degli utenti all'interno del kernel; utilizzava systemcall, anch'esse scritte appositamente per TCFS, per richiedere e salvare le chiavi degli utenti. Inoltre il FS era capace di comunicare solo con NFS, difatti era una versione modificata proprio di NFS.

\subsection{I limiti e le critiche al vecchio TCFS} 
Nonostante TCFS fosse un passo avanti per l'epoca ad oggi appaiono palesi alcune limitazioni date dall'avanzare dei tempi. Con questo si intende che TCFS era limitato alla comunicazione tramite NFS, richiedeva la ri-compilazione del kernel per aggiungere la gestione delle chiavi e richiedeva che le funzioni crittografiche fossero scritte da zero ed internamente al kernel, tutte queste limitazioni sono state superate dalla versione attuale di TCFS. Per quanto riguarda le critiche alla versione originale di TCFS abbiamo che "\textit{TCFS presenta diversi punti deboli che lo rendono meno utile per la distribuzione. In primo luogo, il ricorso alle password di accesso come chiavi utente non è sicuro. Inoltre, la memorizzazione delle chiavi di crittografia su disco in un database di chiavi riduce ulteriormente la sicurezza. Infine, TCFS è disponibile solo su sistemi con kernel Linux 2.2.17 o precedenti, il che ne limita la disponibilità.}"\cite{AnalisiDiDiversiFS}


\section{Il nuovo TCFS: innovazione e sviluppo} 
La nuova versione di TCFS basa la sua implementazione su FUSE, che ha permesso di portare il codice in userspace, utilizzare librerie crittografiche (nel nostro caso fornite da openssl), e di svincolarci da NFS. Di fatti la nuova versione può funzionare insieme a qualsiasi tipo di FS a patto che supporti gli extended attributes. Nella pratica possiamo immaginarlo come un overlay posto al di sopra della cartella sulla quale viene montato TCFS. Lo sviluppo è stato al dir poco travagliato e ha richiesto la costante e continua ricerca delle nuove tecnologie e di come queste potessero essere adattate all'idea proposta. La nuova gestione del keyring di Linux ha permesso di eliminare la necessità di un modulo kernel per la gestione delle chiavi utente, FUSE ha permesso di portare il codice in userspace, MariaDB e GO hanno permesso un implementazzione del threshold sharing semplificata. Il progetto è stato totalmente reinventato.

\subsection{Le ragioni e gli obiettivi della riscrittura del progetto} 
Come specificato precedentemente il codice originale di TCFS è stato perso, per questo motivo si è dovuto riscrivere interamente il FS. Chiaramente dato l'avanzare dei tempi si è ritenuto necessario indagare le nuove tecnologie proposte da quest'epoca, ma si è anche voluto rispettare quanto più possibile il disegno originale, tenendo a mente i seguenti obbiettivi:
\begin{itemize}
    \item \textbf{Portabilità}: Tutte le tecnologie utilizzate da TCFS sono disponibili per OS UNIX-based, cioè Linux, MacOS, NetBSD, OpenBSD, FreeBSD e similari, ma anche per WINDOWS e Solaris, ed è stato scritto in linguaggi portabili, C e GO, rendendo il porting verso nuove architetture plausibile.
    \item \textbf{Semplicità}: TCFS vuole essere semplice nella sua architettura, semplice nel debug e semplice da utilizzare per l'utente.
    \item \textbf{Trasparenza}: La trasparenza in un FS è un concetto fondamentale che richiede una comprensione approfondita. Quando parliamo di trasparenza in un FS, ci riferiamo alla capacità del sistema di nascondere la complessità delle sue operazioni interne all’utente e in questo caso anche al server remoto, rendendo l’interazione con il sistema il più semplice e intuitiva possibile. La creazione di un FS che sia totalmente trasparente all’utente non è un compito semplice per diversi motivi: Le operazioni che un FS esegue in background sono estremamente complesse. Queste includono la gestione dello spazio su disco, l’organizzazione dei file, la gestione dei permessi di accesso, e numerose altre. Nascondere queste operazioni all’utente richiede un design accurato e un’implementazione meticolosa. Quando un FS interagisce con un server remoto, ci sono molte variabili da considerare. Queste includono la latenza della rete, la gestione degli errori di rete, la sincronizzazione dei file tra il client e il server, ecc. Decidere di non applicare modifiche a un server remoto può sembrare una soluzione facile, ma in realtà, raggiungere una trasparenza totale sia verso il server che verso l’utente richiede di tenere conto di queste variabili. Un altro aspetto importante della trasparenza è la gestione degli errori. Un FS deve essere in grado di gestire gli errori in modo tale da minimizzare l’impatto sull’utente. Questo può includere il recupero da errori di sistema, la gestione di errori di rete, la prevenzione della perdita di dati, ecc. Infine, un FS deve essere in grado di fornire un alto livello di performance pur mantenendo la trasparenza. Questo significa che deve essere in grado di gestire grandi quantità di dati, eseguire operazioni in modo efficiente e rispondere rapidamente alle richieste dell’utente. In conclusione, la trasparenza in un FS è un obiettivo difficile da raggiungere a causa della complessità delle operazioni coinvolte e delle sfide associate all’interazione con i server remoti. Tuttavia, con un design accurato e un’implementazione meticolosa, è possibile creare un FS che fornisca un alto livello di trasparenza all’utente.
    \item \textbf{Sostenibilità}: La sostenibilità è un aspetto cruciale nella scrittura di un progetto software. Essa implica la capacità del progetto di resistere nel tempo, adattandosi alle nuove tecnologie e alle esigenze mutevoli degli utenti. La sostenibilità di un progetto software può essere influenzata da vari fattori, tra cui l’uso accurato delle tecnologie disponibili, la manutenibilità del codice e la capacità di adattarsi alle nuove versioni delle dipendenze. Nel caso specifico di TCFS, la versione originale è stata sviluppata come modulo del kernel. Questo significa che ogni volta che viene rilasciata una nuova versione del kernel, è necessario rivedere e potenzialmente aggiornare l’implementazione di TCFS. Questo può rendere la manutenzione del progetto difficile e laboriosa. Tuttavia, la nuova versione di TCFS supera queste limitazioni. Invece di richiedere che il team di TCFS mantenga il modulo del kernel, questa responsabilità viene delegata agli sviluppatori del kernel stesso. Questo significa che, presumibilmente, il modulo del kernel sarà supportato per tutta la sua vita, indipendentemente dalle modifiche al kernel. Questo approccio ha diversi vantaggi. Prima di tutto, riduce il carico di lavoro del team di TCFS, permettendo loro di concentrarsi su altre aree del progetto, non aumentando il carico sugli sviluppatori del kernel. In conclusione, la sostenibilità di un progetto software è un obiettivo importante che può essere raggiunto attraverso un uso accurato delle tecnologie disponibili e un’attenta pianificazione della manutenzione e dell’evoluzione del progetto. Nel caso di TCFS, la delega della manutenzione del modulo del kernel agli sviluppatori del kernel stesso è un esempio di come questo obiettivo possa essere raggiunto a costo zero.
\end{itemize}

\subsection{Le prospettive future del nuovo TCFS}
Al momento di redazione di questa tesi l'implementazione di TCFS risulta incompleta, di fatti non supporta nessun meccanismo di key-recovery e l'implementazione del threshold sharing è parziale, inoltre non è stato svilupparo nessun programma che aiuti l'utente a montare il FS; un altro passo importante è l'utilizzo di funzioni di lettura e scrittura crittografiche che supportino il parallelismo. Ci si prospetta l'implementazione di queste funzionalità nelle versione successive di TCFS. 

\subsection{Il contributo e il valore del nuovo TCFS} 
L’implementazione del nuovo TCFS si colloca in un contesto storico-sociale in cui la salvaguardia dei dati personali è emersa come un problema di rilevanza critica. Questa problematica è particolarmente accentuata nell’attuale era digitale, caratterizzata dalla presenza di conglomerati tecnologici che acquisiscono e analizzano dati personali per scopi di lucro. In questo scenario, la privacy appare come un concetto che rischia di essere trascurato. Come emerge da diverse fonti \cite{4clegal} \cite{altalex1} \cite{altalex2} \cite{cybersecurity360}, le Big Tech sono in grado, attraverso l’utilizzo di strumenti e risorse, capaci di ricavare valore dai dati in loro possesso, di offrire i nostri dati agli advertisers. Tuttavia, questo comporta anche dei rischi e delle sfide, sia per gli utenti che per le autorità, in termini di: diritti e libertà fondamentali della persona; sicurezza e integrità dei dati; concorrenza e mercato; etica e responsabilità.
In questo scenario, il nuovo TCFS offre agli utenti la possibilità di avere il controllo sui propri dati, e di decidere chi e come può accedervi, senza dover rinunciare alla comodità e alla trasparenza di un filesystem tradizionale. In questo modo, il nuovo TCFS contribuisce a rafforzare il principio di protezione delle persone fisiche con riguardo al trattamento dei dati di carattere personale, sancito dal GDPR, e a promuovere una maggiore consapevolezza e gestione dei rischi e della sicurezza legati ai dati.

Uno degli aspetti più importanti e distintivi del nuovo TCFS è la sua capacità di criptare i dati prima di inviarli su un server remoto, tramite protocolli di rete. In questo modo, TCFS garantisce che i dati siano protetti sia in transito che in riposo, cioè sia durante la comunicazione tra il client e il server, sia quando sono memorizzati sul server. Questa funzionalità rende TCFS adatto per l’uso di filesystem di rete, che offrono dei vantaggi, come la condivisione, la scalabilità, e la disponibilità dei dati, ma presentano anche dei rischi, come la perdita, il furto, o la corruzione dei dati, dovuti a possibili attacchi o malfunzionamenti della rete. Con il nuovo TCFS, gli utenti possono usufruire dei servizi offerti dalle grandi aziende tecnologiche, senza dover esporre i propri dati in chiaro, e senza dover affidare la loro confidenzialità a terze parti. In questo modo, il nuovo TCFS tutela la privacy degli utenti, e li rende più indipendenti e sicuri.

\subsection{Le considerazioni personali e le lezioni apprese}
Nell’ambito dell’attuale ecosistema tecnologico, l’insorgere di progetti come TCFS appare non solo appropriato, ma imprescindibile. L’appeal di servizi di hosting gratuiti è indiscutibilmente rilevante, tuttavia, è essenziale riconoscere che in tali contesti, l’utente si trasforma inevitabilmente nel prodotto. È altrettanto rilevante sottolineare che, nonostante l’esistenza di servizi di hosting a pagamento, la riservatezza dei dati non è assicurata in quasi nessuna circostanza. Pertanto, l’integrazione di TCFS in questo scenario si rivela di fondamentale importanza. Grazie alla sua eccezionale facilità d’uso, TCFS può essere adoperato anche da utenti con competenze limitate, offrendo una soluzione semplice ma efficace alle sfide inerenti alla privacy che caratterizzano l’era digitale contemporanea.

L’implementazione di questo progetto, come precedentemente accennato, ha comportato notevoli sfide. Questo è dovuto alla vastità del progetto e alla ricerca delle tecnologie appropriate, che ha richiesto un impegno temporale considerevole. Inoltre, l’implementazione degli algoritmi crittografici di OpenSSL si è rivelata un compito arduo. L’acquisizione della padronanza di tutte le tecnologie coinvolte è stata un’esperienza estremamente arricchente. Concludo questo progetto con una consapevolezza più profonda della complessità di AES, delle sfide relative alla privacy, del kernel Linux e dei linguaggi di programmazione C, Go e SQL.
